%قبل از شروع مسئله ترجیحاً چیزی ننویسید. اما اگر خواستید بنویسید، در قالب لیستی (فقرات و شمارش) بنویسید.
%برای اضافه کردن نمره، بعدی از شروع‌مسئله، [x] را اضافه کنید و به جای x  نمره مورد نظر را قرار دهید. برای زیرمسئله‌ها هم به همین صورت است.
%\شروع‌مسئله[10]
\شروع‌مسئله
در جدول درهم سازی با استفاده از روش وارسی خطی، تابع درهم سازی برای جدولی با اندازه هشت به صورت زیر است:


\begin{center}
	\begin{tabular}{|c|c|c|c|c|c|c|c|c| }
		\hline
		H & G & F & E & D & C & B & A & key\\ \hline
		1 & 4 & 5 & 4 & 4 & 2 & 6 & 2 & hash \\
		\hline
	\end{tabular}
\end{center}

اگر جدول درهم سازی در ابتدا تهی باشد، به چند حالت میتوان این عناصر را در جدول درج کرد تا در نهایت جدول زیر تولید شود؟

\begin{center}
	\begin{tabular}{|c|c|c|c|c|c|c|c|c| }
		\hline
		7 & 6 & 5 & 4 & 3 & 2 & 1 & A & i\\ \hline
		F & B & E & D & C & A & H & G & T(i) \\
		\hline
	\end{tabular}
\end{center}


\solution{
\begin{flushleft}
	$Hash(A) = 2, Hash(C) = 2$\\
	$T(2) = A, T(3) = C$\\
\end{flushleft}
	بنابراین به دو مورد زیر C بعد از A هش میشود.
\begin{flushleft}
	$Hash(D) = 4, Hash(E) = 4$\\
	$T(4) = D, T(5) = E$\\
\end{flushleft}
و E بعد از D هش میشود.

\begin{flushleft}
	$Hash(E) = 4, Hash(B) = 6, Hash(F) = 5$\\
	$T(5) = E, T(6) = B, T(7) = F$\\
\end{flushleft}
و F بعد از E و B هش میشود.

\begin{flushleft}
	$Hash(D) = 4, Hash(E) = 4, Hash(G) = 4$\\
	$Hash(B) = 6, Hash(F) = 5$\\
	$T(5) = E, T(6) = B$\\
	$T(7) = F, T(0) = G$\\
\end{flushleft}
 و G باید بعد از F و E و B هش شود.
 
 حالا تعداد جایشگت های مختلف را محاسبه میکنیم. H هیج محدودیتی نداشته است پس 8 جایگشت مختلف دارد. A و C هم با توجه به یکی از محدودیت ها دارای $C(7, 2)$ یعنی 21 جایگشت مختلف هستند. با توجه به شروط دو و سه و چهار هم متوجه میشویم که ترتیب هش شدن به صورت D و E و F و G است و B هم قبل از F باید هش شود پس 3 حالت دارد.  
 
 برابراین به طور کلی $8 * 3* 21 = 504$ حالت مختلف برای ساختن این آرایه داریم.
 
 
}